% AA-03 | Algebra abstrakcyjna | kanoniczny rozdzial: dzialania grup
% Zachowano istniejace etykiety, twierdzenia, dowody i przyklady.
\chapter{Działania grup, orbity i równanie klas}
\label{aa:ch:dzialania}

\section{Dlaczego grupy działają na zbiorach?}
Grupa koduje działania odwracalne, lecz sama informacja o jej mnożeniu
nie mówi jeszcze, \emph{na czym} wykonujemy te operacje. Grupa
permutacji przestawia elementy zbioru, obroty przesuwają punkty sfery,
a macierze odwracalne przekształcają wektory. Wspólną definicję
otrzymamy, gdy narzucimy zgodność kolejnych operacji z mnożeniem w grupie.

\begin{definicja}[Lewe działanie grupy]
\label{aa:def:dzialanie}
Niech $G$ będzie grupą z elementem neutralnym $e$, a $X$ zbiorem.
\emph{Lewym działaniem} $G$ na $X$ nazywamy odwzorowanie
$G\times X\to X$, $(g,x)\mapsto g\cdot x$, spełniające
$e\cdot x=x$ i $(gh)\cdot x=g\cdot(h\cdot x)$
dla wszystkich $g,h\in G$ i $x\in X$.
Zbiór $X$ z takim działaniem nazywamy \emph{$G$-zbiorem}.
Punkt $x\in X$ jest \emph{stały} dla działania całej grupy,
jeżeli $g\cdot x=x$ dla wszystkich $g\in G$.
\end{definicja}

\begin{lemat}[Działanie to homomorfizm do grupy permutacji]
\label{aa:lem:akcja-hom}
Dane jest lewe działanie grupy $G$ na zbiorze $X$.
Przez $\operatorname{Sym}(X)$ oznaczamy grupę wszystkich bijekcji
$X\to X$ z działaniem składania odwzorowań. Dla każdego
$g\in G$ odwzorowanie $T_g:X\to X$, $T_g(x)=g\cdot x$,
jest bijekcją, a funkcja
$\rho:G\to\operatorname{Sym}(X)$, $\rho(g)=T_g$,
jest homomorfizmem grup.
Odwrotnie, każdy homomorfizm $\rho:G\to\operatorname{Sym}(X)$
wyznacza lewe działanie przez $g\cdot x=\rho(g)(x)$.
\end{lemat}
\begin{proof}
Z aksjomatu działania $T_e(x)=e\cdot x=x$, więc $T_e=\operatorname{id}_X$.
Ponadto $T_g(T_{g^{-1}}(x))=(gg^{-1})\cdot x=x$;
po zamianie $g$ z $g^{-1}$ dostajemy także
$T_{g^{-1}}(T_g(x))=x$. Stąd $T_g^{-1}=T_{g^{-1}}$,
a więc $T_g$ jest bijekcją.
Dla wszystkich $g,h,x$ mamy
$T_{gh}(x)=(gh)\cdot x=g\cdot(h\cdot x)=(T_g\circ T_h)(x)$.
Równość dla każdego $x$ dowodzi $\rho(gh)=\rho(g)\rho(h)$.
W drugą stronę homomorfizm zachowuje element neutralny,
więc $\rho(e)=\operatorname{id}_X$, a równość
$\rho(gh)=\rho(g)\circ\rho(h)$ daje oba aksjomaty działania.
\end{proof}

\begin{definicja}[Działanie wierne i przechodnie]
Niech grupa $G$ działa na zbiorze $X$, a
$\rho:G\to\operatorname{Sym}(X)$ będzie homomorfizmem
odpowiadającym temu działaniu. Działanie jest \emph{wierne}, jeżeli
$\rho$ jest różnowartościowy;
jest \emph{przechodnie}, jeżeli dla wszystkich $x,y\in X$
istnieje $g\in G$ z $g\cdot x=y$.
Jądro działania to $\ker\rho$.
\end{definicja}

\begin{przyklad}[Trzy działania jednej grupy]
\label{aa:prz:trzy-akcje}
Niech $G=S_3$.
\begin{enumerate}
\item $S_3$ działa na $\{1,2,3\}$ przez $\sigma\cdot i=\sigma(i)$.
Jest to działanie wierne i przechodnie: każdą liczbę można
przenieść na dowolną inną, a permutacja ustalająca wszystkie
trzy liczby jest identycznością.
\item $S_3$ działa na zbiorze własnych elementów przez mnożenie
z lewej: $\sigma\cdot\tau=\sigma\tau$. To także działanie wierne
i przechodnie: aby przenieść $\tau$ na $\eta$, bierzemy
$\sigma=\eta\tau^{-1}$. Działanie to ma sześciopunktową,
a nie trzypunktową przestrzeń.
\item $S_3$ działa na sobie przez sprzężenie
$\sigma\cdot\tau=\sigma\tau\sigma^{-1}$.
Nie jest przechodnie: element neutralny przechodzi tylko na siebie.
Każda permutacja w centrum działa trywialnie; pokażemy,
że właśnie centrum stanowi jądro tego działania.
\end{enumerate}
To, że grupy są te same, nie oznacza, że ich działania są takie same.
\end{przyklad}

\begin{przyklad}[Działanie przechodnie nie musi być wierne]
\label{aa:prz:akcja-c4-dwa-punkty}
Niech $G=\mathbb Z/4\mathbb Z$, a $X=\{u,v\}$.
Dla klasy $[k]\in G$ określamy $[k]\cdot u=u$ i
$[k]\cdot v=v$, jeśli $k$ jest parzyste, oraz
$[k]\cdot u=v$, $[k]\cdot v=u$, jeśli $k$ jest nieparzyste.
Parzystość nie zależy od reprezentanta klasy modulo $4$.
Zgodność działania z dodawaniem wynika stąd, że złożenie
zamiany punktów ze sobą daje identyczność, a parzystość sumy
jest sumą parzystości modulo $2$.
Zatem $[0]$ działa identycznościowo, a
$([k]+[\ell])\cdot x=[k]\cdot([\ell]\cdot x)$ dla
$k,\ell\in\mathbb Z$ i $x\in X$.
Działanie jest przechodnie: $[1]\cdot u=v$ i $[1]\cdot v=u$,
a $[0]$ ustala oba punkty.
Nie jest wierne: $[2]\ne[0]$ ustala oba punkty,
a dokładnie klasy parzyste działają trywialnie,
więc $\ker\rho=\{[0],[2]\}$.
Dokładnie dwie klasy $[0]$ i $[2]$ ustalają punkt $u$,
a jego możliwymi obrazami są oba punkty $u,v$.
Zależność $2\cdot2=4$ wyjaśnimy ogólnie w twierdzeniu
o orbicie i stabilizatorze. Przykład pokazuje, że
przechodniość nie pociąga wierności.

\begin{center}
\begin{tikzpicture}[>=Latex,baseline=(current bounding box.center)]
\node[draw,circle,minimum size=9mm] (u) at (0,0) {$u$};
\node[draw,circle,minimum size=9mm] (v) at (3,0) {$v$};
\draw[->,bend left=24] (u) to node[above] {$[1]$} (v);
\draw[->,bend left=24] (v) to node[below] {$[1]$} (u);
\node[align=center] at (1.5,-1.55) {$[2]$ ustala oba punkty};
\end{tikzpicture}
\end{center}
\end{przyklad}

\section{Orbity i stabilizatory}
\begin{definicja}[Orbita i stabilizator]
\label{aa:def:orbita-stabilizator}
Niech $G\curvearrowright X$ będzie działaniem, a $x\in X$.
\emph{Orbitą} punktu $x$ nazywamy zbiór
$G\cdot x=\{g\cdot x:g\in G\}$.
\emph{Stabilizatorem} punktu $x$ nazywamy
$G_x=\{g\in G:g\cdot x=x\}$.
\end{definicja}

\begin{lemat}[Stabilizator jest podgrupą, orbity rozkładają zbiór]
\label{aa:lem:orbity-podzial}
Niech $G$ działa na zbiorze $X$.
Dla każdego $x\in X$ zbiór $G_x$ jest podgrupą $G$.
Relacja $x\sim y\iff y\in G\cdot x$ jest relacją równoważności,
a jej klasy to dokładnie orbity działania.
\end{lemat}
\begin{proof}
Neutralny element $e$ ustala $x$, więc $e\in G_x$.
Jeżeli $g,h\in G_x$, to
$(gh^{-1})\cdot x=g\cdot(h^{-1}\cdot x)=g\cdot x=x$:
istotnie, skoro $h\cdot x=x$, zastosowanie $h^{-1}$
po obu stronach daje $h^{-1}\cdot x=x$.
Kryterium podgrupy z Rozdziału~\ref{aa:ch:grupy} daje $G_x\le G$.

Relacja jest zwrotna, bo $x=e\cdot x$.
Jeśli $y=g\cdot x$, to $x=g^{-1}\cdot y$, zatem jest symetryczna.
Jeśli ponadto $z=h\cdot y$, to $z=(hg)\cdot x$, zatem
jest przechodnia. Jej klasą równoważności punktu $x$
jest zbiór wszystkich $g\cdot x$, czyli $G\cdot x$.
Klasy równoważności są albo równe, albo rozłączne
i pokrywają $X$.
\end{proof}

\begin{twierdzenie}[Orbita i stabilizator]
\label{aa:tw:orbita-stabilizator}
Niech grupa $G$ działa na zbiorze $X$, a $x\in X$.
Odwzorowanie
$\Theta_x:G/G_x\to G\cdot x$, $\Theta_x(gG_x)=g\cdot x$,
jest dobrze określoną bijekcją zbiorów. W szczególności,
gdy $G$ jest skończona, zachodzi
\[
 |G\cdot x|=[G:G_x]=\frac{|G|}{|G_x|}.
\]
\end{twierdzenie}
\begin{proof}
Najpierw porównajmy dwa reprezentanty warstwy:
$gG_x=hG_x$ dokładnie wtedy, gdy $h^{-1}g\in G_x$,
czyli $(h^{-1}g)\cdot x=x$.
Z aksjomatów działania ostatnia równość jest równoważna
$g\cdot x=h\cdot x$: zastosowanie $h$ daje implikację
w prawo, a zastosowanie $h^{-1}$ daje odwrotną.
Wykazaliśmy jednocześnie, że wartość $\Theta_x$
nie zależy od reprezentanta warstwy i że równe obrazy
pochodzą tylko z równych warstw; jest więc różnowartościowa.
Każdy element orbity ma postać $g\cdot x$ i jest obrazem
warstwy $gG_x$, zatem $\Theta_x$ jest surjektywna.
Dla skończonego $G$ stosujemy twierdzenie Lagrange'a
(\ref{aa:tw:lagrange}) do $G_x\le G$.
\end{proof}

\begin{przyklad}[Naturalne działanie $S_4$]
Niech $S_4$ działa na $\{1,2,3,4\}$ w zwykły sposób.
Orbita punktu $1$ to cały zbiór czteroelementowy:
permutacją można wysłać $1$ na dowolną liczbę $i$.
Stabilizator $1$ składa się z permutacji pozostałych trzech
liczb i ma $3!=6$ elementów. Twierdzenie daje
$|S_4\cdot1|=|S_4|/|(S_4)_1|=24/6=4$.
Widać, co dzieje się w dowodzie: sześć permutacji
wysyłających $1$ na określoną liczbę tworzy jedną warstwę
stabilizatora; cztery możliwe obrazy to cztery takie warstwy.
\end{przyklad}

\begin{przyklad}[Nieskończona grupa z orbitą skończoną]
\label{aa:prz:z-dziala-na-zn}
Niech $n\ge2$ będzie liczbą naturalną.
Grupa $(\mathbb Z,+)$ działa na $X=\mathbb Z/n\mathbb Z$
przez $k\cdot[a]=[a+k]$.
Jeśli $[a]=[b]$, to $n$ dzieli $a-b$, a zatem dzieli
$(a+k)-(b+k)$, więc działanie jest dobrze określone.
Ponadto $0\cdot[a]=[a]$ i
$(k+\ell)\cdot[a]=[a+k+\ell]=k\cdot(\ell\cdot[a])$.
Jest przechodnie: aby przenieść $[a]$ na $[b]$, wystarczy
wybrać $k=b-a$.
Stabilizator $[0]$ to $n\mathbb Z$, gdyż
$k\cdot[0]=[0]$ dokładnie wtedy, gdy $n\mid k$.
Warstwy $\mathbb Z/n\mathbb Z$ grupy addytywnej
$\mathbb Z$ względem $n\mathbb Z$ są dokładnie $n$
klasami reszt, więc
$|\mathbb Z\cdot[0]|=[\mathbb Z:n\mathbb Z]=n$.
Nie stosujemy tutaj ilorazu $|\mathbb Z|/|n\mathbb Z|$:
oba zbiory są nieskończone, natomiast ich indeks jest skończony.
\end{przyklad}

\begin{przyklad}[Działanie na klasach warstw]
\label{aa:prz:akcja-kosety}
Niech $H\le G$. Grupa $G$ działa na zbiorze $G/H$
przez $g\cdot(aH)=(ga)H$.
Wzór jest dobrze określony: jeśli $aH=bH$, to
$b^{-1}a\in H$, a więc
$(gb)^{-1}(ga)=b^{-1}a\in H$, czyli $(ga)H=(gb)H$.
Aksjomaty działania wynikają z
$e\cdot(aH)=aH$ oraz
$(gh)\cdot(aH)=((gh)a)H=g\cdot((ha)H)
=g\cdot(h\cdot(aH))$.
Jest to działanie przechodnie, ponieważ
$(ba^{-1})\cdot(aH)=bH$.
Stabilizator warstwy $H$ to dokładnie $H$, gdyż
$gH=H\iff g\in H$.
Stabilizator warstwy $aH$ wynosi $aHa^{-1}$:
$g\cdot aH=aH$ wtedy i tylko wtedy, gdy
$a^{-1}ga\in H$.
To ważny model wszystkich działań przechodnich.
\end{przyklad}

\begin{wniosek}[Każdy przechodni $G$-zbiór jest przestrzenią warstw]
\label{aa:wn:przechodnie-warstwy}
Niech grupa $G$ działa przechodnio na niepustym zbiorze $X$.
Dla dowolnego $x\in X$ istnieje bijekcja
$G/G_x\to X$, $gG_x\mapsto g\cdot x$,
która jest zgodna z działaniem $G$ po obu stronach.
\end{wniosek}
\begin{proof}
Z przechodniości $G\cdot x=X$, więc bijekcja wynika
z twierdzenia o orbicie i stabilizatorze.
Dla $a,g\in G$ zgodność działania oznacza
$\Theta_x(a\cdot(gG_x))=\Theta_x((ag)G_x)
=(ag)\cdot x=a\cdot(g\cdot x)=a\cdot\Theta_x(gG_x)$.
\end{proof}

\section{Jądro działania i normalność}
\begin{twierdzenie}[Jądro i efektywne działanie]
\label{aa:tw:jadro-akcji}
Jeżeli grupa $G$ działa na zbiorze $X$ przez homomorfizm
$\rho:G\to\operatorname{Sym}(X)$, to
\[
 \ker\rho=\{g\in G:g\cdot x=x\text{ dla każdego }x\in X\}
 =\bigcap_{x\in X}G_x.
\]
Zbiór ten jest podgrupą normalną, a działanie jest wierne
dokładnie wtedy, gdy to jądro jest trywialne.
Działanie przechodzi na wierne działanie $G/\ker\rho$
na tym samym zbiorze $X$.
\end{twierdzenie}
\begin{proof}
Z definicji $g\in\ker\rho$ oznacza
$\rho(g)=\operatorname{id}_X$, czyli
$\rho(g)(x)=x$ dla każdego $x\in X$.
To jest równoważne przynależności do wszystkich stabilizatorów.
Jądro homomorfizmu jest normalne, a homomorfizm jest
różnowartościowy dokładnie wtedy, gdy jego jądro jest trywialne
(Rozdział~\ref{aa:ch:lagrange-ilorazy}).
Niech $N=\ker\rho$. Definiujemy
$(gN)\cdot x=g\cdot x$.
Jeśli $gN=hN$, to $h^{-1}g\in N$;
więc $(h^{-1}g)\cdot x=x$, co po zastosowaniu $h$
daje $g\cdot x=h\cdot x$. Działanie jest dobrze określone. Ponadto
$(eN)\cdot x=e\cdot x=x$ oraz
$((gN)(hN))\cdot x=(gh)\cdot x
=g\cdot(h\cdot x)=(gN)\cdot((hN)\cdot x)$.
Jeśli warstwa $gN$ ustala wszystkie $x$, to $g\in N$
i $gN=N$, więc nowe działanie jest wierne.
\end{proof}

\begin{przyklad}[Działanie na warstwach i rdzeń podgrupy]
Niech $H\le G$ i rozważmy działanie na $G/H$.
Element $g$ działa trywialnie wtedy i tylko wtedy, gdy
$g$ ustala wszystkie warstwy $aH$.
Stabilizator $aH$ to $aHa^{-1}$, stąd jądro działania
wynosi $\bigcap_{a\in G}aHa^{-1}$.
Jest to \emph{rdzeń normalny} podgrupy $H$, czyli
największa normalna podgrupa $G$ zawarta w $H$.
Rzeczywiście jest normalna jako jądro i zawiera się
w $H$ (bierzemy $a=e$).
Jeżeli $N\trianglelefteq G$ oraz $N\subseteq H$, to
$aNa^{-1}=N\subseteq aHa^{-1}$ dla każdego $a$,
a więc $N$ zawiera się w tym przecięciu.
\end{przyklad}

\begin{przyklad}[Wierne działanie na warstwach nie wymaga normalności]
\label{aa:prz:s3-core}
W $G=S_3$ niech $H=\langle(2\,3)\rangle$, czyli
stabilizator punktu $1$ w naturalnym działaniu na
$\{1,2,3\}$. Zbiór warstw $G/H$ ma trzy elementy.
Odwzorowanie $gH\mapsto g(1)$ jest bijekcją zgodną
z działaniem $G$; wynika to z Twierdzenia~\ref{aa:tw:orbita-stabilizator}.
Jądro działania na trzech warstwach jest trywialne, bo
permutacja ustalająca jednocześnie punkty $1,2,3$
jest identycznością. Zatem rdzeń normalny $H$ jest
równy $\{e\}$. Sama $H$ nie jest jednak normalna:
$(1\,2)(2\,3)(1\,2)=(1\,3)\notin H$.
W szczególności zbiór $G/H$ ma naturalne działanie $G$,
lecz nie ma dobrze określonego mnożenia warstw
$(gH)(hH)=(gh)H$, które tworzyłoby grupę ilorazową.
Rozróżnienie to wykorzystamy w następnym rozdziale,
w dowodzie drugiego twierdzenia Sylowa.
\end{przyklad}

\section{Sprzężenia, centralizatory i normalizatory}
\begin{definicja}[Centralizator i normalizator]
\label{aa:def:centralizator-normalizator}
Niech $G$ będzie grupą, $x\in G$, $H\le G$.
\emph{Centralizator} $x$ to
$C_G(x)=\{g\in G:gx=xg\}$, a \emph{normalizator} $H$ to
$N_G(H)=\{g\in G:gHg^{-1}=H\}$.
\emph{Klasa sprzężoności} $x$ to
$\operatorname{Cl}_G(x)=\{gxg^{-1}:g\in G\}$.
\end{definicja}

\begin{lemat}[Centralizatory i normalizatory jako stabilizatory]
\label{aa:lem:centralizator-normalizator}
Dla dowolnej grupy $G$, elementu $x\in G$ i podgrupy $H\le G$
zbiory $C_G(x)$ i $N_G(H)$ są podgrupami $G$.
Ponadto $H\trianglelefteq N_G(H)$ oraz
\[
 |\operatorname{Cl}_G(x)|=[G:C_G(x)]
\]
przy skończonym $G$.
\end{lemat}
\begin{proof}
Rozważmy działanie $G$ na sobie przez sprzężenie:
$g\cdot x=gxg^{-1}$.
Aksjomaty działania wynikają z
$e x e^{-1}=x$ oraz
$(gh)x(gh)^{-1}=g(hxh^{-1})g^{-1}$.
Stabilizatorem $x$ jest zbiór tych $g$, dla których
$gxg^{-1}=x$, czyli $gx=xg$.
Zatem $C_G(x)=G_x$ i jest podgrupą, a formuła
na liczebność klasy sprzężoności wynika z twierdzenia
o orbicie i stabilizatorze.

Grupa $G$ działa również na zbiorze wszystkich podgrup $G$
przez $g\cdot K=gKg^{-1}$.
Aksjomaty działania sprawdza się jak wyżej,
a sprzężona podgrupa jest podgrupą dzięki zachowaniu
mnożenia i odwrotności przez sprzężenie.
Stabilizatorem podgrupy $H$ jest $N_G(H)$, stąd jest podgrupą.
Dla każdego $h\in H$ zachodzi $hHh^{-1}=H$,
więc $H\le N_G(H)$. Ponadto każdy element $N_G(H)$
sprzęga $H$ na $H$, czyli $H\trianglelefteq N_G(H)$.
\end{proof}

\begin{przyklad}[Klasy sprzężoności w $S_3$]
W $S_3$ są: identyczność, trzy transpozycje i dwa $3$-cykle.
Przez sprzężenie permutacje zachowują strukturę cykli:
dla bijekcji $g$ zachodzi
$g(a_1\,\ldots\,a_k)g^{-1}=(g(a_1)\,\ldots\,g(a_k))$,
co sprawdzamy, stosując obie strony do $g(a_i)$;
poza punktami $g(a_1),\ldots,g(a_k)$ obie strony
ustalają pozostałe punkty. Dla dowolnych dwóch cykli
tej samej długości możemy wybrać permutację $g$, która
przenosi kolejne elementy pierwszego cyklu na kolejne
elementy drugiego. Stąd trzy transpozycje są sprzężone,
a dwa $3$-cykle również są sprzężone.
Centralizator transpozycji $(1\,2)$ zawiera tylko
$e$ i $(1\,2)$: jeśli $g(1\,2)g^{-1}=(1\,2)$,
to $g$ musi zachować dwuelementowy zbiór $\{1,2\}$,
co pozostawia dwa wybory. Ma więc rząd $2$ i klasa
ma $6/2=3$ elementy.
Centralizator $(1\,2\,3)$ to jego trzy potęgi:
możemy wysłać $1$ na jeden z trzech punktów,
a komutowanie z cyklem wymusza obrazy $2$ i $3$.
Zatem klasa $3$-cykli ma $6/3=2$ elementy.
Ostatecznie klasy sprzężoności mają rozmiary $1,3,2$.
\end{przyklad}

\begin{przyklad}[Jądro działania przez sprzężenie]
Dla dowolnej grupy $G$ jądro działania przez sprzężenie
składa się z elementów $g$ spełniających
$gxg^{-1}=x$ dla wszystkich $x\in G$, czyli z $Z(G)$.
Zatem $G/Z(G)$ działa wiernie na $G$ przez sprzężenie.
Przez $\operatorname{Aut}(G)$ oznaczamy grupę
wszystkich automorfizmów $G$ ze składaniem odwzorowań.
Otrzymujemy przy tym homomorfizm
$G/Z(G)\to\operatorname{Aut}(G)$ dany wzorem
$gZ(G)\mapsto(x\mapsto gxg^{-1})$. Jest on dobrze określony
na mocy opisu jądra działania. Sprzężenie przez $g$
zachowuje mnożenie, a jego odwrotnością jest sprzężenie
przez $g^{-1}$, więc rzeczywiście jest automorfizmem $G$.
Zgodność z mnożeniem warstw wynika ze wzoru
$(gh)x(gh)^{-1}=g(hxh^{-1})g^{-1}$.
\end{przyklad}

\section{Równanie klas i zastosowania do \texorpdfstring{$p$}{p}-grup}
\begin{twierdzenie}[Równanie klas]
\label{aa:tw:rownanie-klas}
Niech $G$ będzie grupą skończoną.
Wybierzmy reprezentantów $x_1,\ldots,x_r$ wszystkich
klas sprzężoności niezawartych w centrum $Z(G)$.
Wówczas
\[
 |G|=|Z(G)|+\sum_{i=1}^{r}[G:C_G(x_i)].
\]
Każdy indeks w występującej sumie jest liczbą całkowitą
nie mniejszą niż $2$.
\end{twierdzenie}
\begin{proof}
Orbity działania przez sprzężenie tworzą rozkład $G$.
Orbita $\operatorname{Cl}_G(x)$ ma jeden element wtedy i tylko wtedy,
gdy $gxg^{-1}=x$ dla wszystkich $g$, co jest równoważne
$x\in Z(G)$. Orbity jednoelementowe składają się zatem
z dokładnie $|Z(G)|$ elementów.
Pozostałe orbity są parami rozłączne i mają rozmiary
$[G:C_G(x_i)]$, ponieważ centralizator jest stabilizatorem.
Sumując liczby elementów wszystkich orbit, otrzymujemy wzór.
Jeśli $x_i\notin Z(G)$, to jego orbita ma więcej niż jeden element.
\end{proof}

\begin{definicja}[Grupa o rzędzie będącym potęgą liczby pierwszej]
Niech $p$ będzie liczbą pierwszą, a $a\ge0$ liczbą całkowitą.
Grupę o rzędzie $p^a$ nazywamy \emph{$p$-grupą}.
Jeżeli $a=0$, jest to grupa trywialna, ponieważ ma tylko
jeden element. W dalszych twierdzeniach o nietrywialnym
centrum zakładamy $a\ge1$.
\end{definicja}

\begin{twierdzenie}[Centrum niebanalnej skończonej $p$-grupy]
\label{aa:tw:centrum-pgrupy}
Jeżeli $G$ jest grupą rzędu $p^a$, gdzie $p$ jest liczbą pierwszą
i $a\ge1$, to $|Z(G)|$ jest podzielne przez $p$.
W szczególności centrum zawiera element różny od $e$.
\end{twierdzenie}
\begin{proof}
Dla każdego $x\in G$ centralizator $C_G(x)$ jest podgrupą,
więc jego rząd jest potęgą $p$ na mocy twierdzenia Lagrange'a.
Jeśli $x\notin Z(G)$, to $C_G(x)$ jest właściwą podgrupą:
rzeczywiście gdyby $C_G(x)=G$, element $x$ komutowałby
ze wszystkimi elementami grupy i należałby do centrum.
A zatem indeks $[G:C_G(x)]$ jest dodatnią potęgą $p$,
i wynosi co najmniej $p$.
Równanie klas daje
$|G|=|Z(G)|+\sum_i [G:C_G(x_i)]$.
Obie liczby $|G|=p^a$ i każdy składnik sumy są podzielne
przez $p$, więc ich różnica $|Z(G)|$ również.
Ponieważ $e\in Z(G)$, centrum nie jest puste;
jeśli jego liczebność jest dodatnią wielokrotnością $p$,
to wynosi co najmniej $p\ge2$.
\end{proof}

\begin{lemat}[Cykliczny iloraz przez centrum wymusza przemienność]
\label{aa:lem:cykliczny-centrum}
Jeśli $G$ jest grupą i $G/Z(G)$ jest cykliczna,
to $G$ jest abelowa.
\end{lemat}
\begin{proof}
Niech $gZ(G)$ generuje $G/Z(G)$.
Dla dowolnych $x,y\in G$ istnieją liczby całkowite $m,n$
i elementy $z,w\in Z(G)$ z $x=g^m z$, $y=g^n w$.
Korzystając z tego, że $z,w$ komutują ze wszystkimi
elementami $G$, otrzymujemy
$xy=g^{m+n}zw=g^{n+m}wz=yx$.
Ponieważ $x,y$ były dowolne, grupa jest abelowa.
\end{proof}

\begin{twierdzenie}[Grupy rzędu $p^2$ są abelowe]
\label{aa:tw:p2-abelowa}
Niech $p$ będzie liczbą pierwszą, a $G$ grupą rzędu $p^2$.
Wtedy $G$ jest abelowa.
\end{twierdzenie}
\begin{proof}
Poprzednie twierdzenie daje $p\mid |Z(G)|$,
a twierdzenie Lagrange'a daje $|Z(G)|\mid p^2$.
Stąd $|Z(G)|$ wynosi $p$ albo $p^2$.
Jeśli $|Z(G)|=p^2$, to $Z(G)=G$ i grupa jest abelowa.
Jeśli $|Z(G)|=p$, to iloraz $G/Z(G)$ ma rząd $p$,
a więc jest cykliczny. Poprzedni lemat ponownie daje
przemienność $G$.
\end{proof}

\begin{przyklad}[Dwie grupy abelowe rzędu cztery]
Grupa $\mathbb Z/4\mathbb Z$ ma element $[1]$ rzędu $4$,
natomiast $(\mathbb Z/2\mathbb Z)^2$ ma trzy niezerowe
elementy, z których każdy ma rząd $2$.
Obie grupy mają rząd $4$ i są abelowe,
lecz nie są izomorficzne: izomorfizm zachowuje rzędy elementów.
Z samej abelowości nie wynika więc cykliczność.
\end{przyklad}

\section{Lemat Burnside'a o orbitach}
\begin{definicja}[Zbiór punktów stałych elementu grupy]
Niech grupa $G$ działa na zbiorze $X$.
Dla $g\in G$ definiujemy
$X^g=\{x\in X:g\cdot x=x\}$.
Nie mylimy $X^g$ ze zbiorem punktów stałych całego działania,
który jest przecięciem wszystkich $X^g$.
\end{definicja}

\begin{twierdzenie}[Lemat Burnside'a]
\label{aa:tw:burnside}
Niech skończona grupa $G$ działa na skończonym zbiorze $X$.
Niech $X/G=\{G\cdot x:x\in X\}$ oznacza zbiór orbit.
Wtedy liczba orbit $|X/G|$ spełnia
\[
 |X/G|=\frac{1}{|G|}\sum_{g\in G}|X^g|.
\]
\end{twierdzenie}
\begin{proof}
Policzymy na dwa sposoby zbiór par
$P=\{(g,x)\in G\times X:g\cdot x=x\}$.
Ustalając najpierw $g$, otrzymujemy
$|P|=\sum_{g\in G}|X^g|$.
Ustalając natomiast $x$, widzimy, że odpowiadające mu
pierwsze współrzędne tworzą stabilizator $G_x$,
a więc $|P|=\sum_{x\in X}|G_x|$.

Podzielmy teraz $X$ na orbity. Jeżeli $O=G\cdot x$
jest jedną z nich, to wszystkie stabilizatory punktów z $O$
mają ten sam rząd: dla $y=h\cdot x$ mamy
$G_y=hG_xh^{-1}$, co wynika z rachunku
$g\cdot(h\cdot x)=h\cdot x\iff
(h^{-1}gh)\cdot x=x$.
Dlatego
$\sum_{y\in O}|G_y|=|O|\,|G_x|$.
Twierdzenie o orbicie i stabilizatorze daje
$|O|\,|G_x|=|G|$.
Każda orbita wnosi zatem dokładnie $|G|$ do sumy
$\sum_{x\in X}|G_x|$.
Jeśli orbit jest $k$, to $|P|=k|G|$.
Porównując oba sposoby liczenia $P$ i dzieląc przez $|G|$,
otrzymujemy żądany wzór.
\end{proof}

\begin{przyklad}[Kolorowanie wierzchołków kwadratu dwoma kolorami]
\label{aa:prz:burnside-kwadrat}
Ustalmy cztery wierzchołki kwadratu i dwa kolory: czarny i biały.
Kolorowanie to funkcja ze zbioru czterech wierzchołków
do zbioru dwuelementowego, więc jest ich $2^4=16$.
Dwa kolorowania uznajemy za to samo, gdy jedno można otrzymać
z drugiego przez symetrię kwadratu. Działa grupa $D_4$
o ośmiu elementach. Jeśli $V$ oznacza zbiór wierzchołków,
a $C=\{\text{czarny},\text{biały}\}$ zbiór kolorów,
to przestrzenią kolorowań jest zbiór funkcji $C^V$.
Dla symetrii $g\in D_4$, kolorowania $c:V\to C$ oraz $v\in V$
definiujemy $(g\cdot c)(v)=c(g^{-1}v)$.
To jest działanie: $e\cdot c=c$, zaś
$((gh)\cdot c)(v)=c(h^{-1}g^{-1}v)
=(g\cdot(h\cdot c))(v)$.
Kolorowanie $c$ jest ustalane przez $g$ wtedy i tylko wtedy,
gdy jest stałe na każdym cyklu permutacji wierzchołków
wyznaczonej przez $g$. Jeżeli cykli jest $k$, można
wybrać jeden z dwóch kolorów niezależnie dla każdego cyklu;
zatem $| (C^V)^g |=2^k$.
Obliczmy dokładnie liczby kolorowań zachowanych przez symetrie.
\begin{enumerate}
\item Identyczność ustala wszystkie $16$ kolorowań.
\item Każdy z dwóch obrotów o $90^\circ$ i $270^\circ$
cyklicznie permutuje wszystkie cztery wierzchołki,
więc kolorowanie jest stałe na całym kwadracie:
każdy z tych obrotów ustala $2$ kolorowania.
\item Obrót o $180^\circ$ wymienia dwie pary wierzchołków,
więc można niezależnie wybrać kolor każdej pary:
ustala $2^2=4$ kolorowania.
\item Każde z dwóch odbić względem przekątnych ustala
dwa wierzchołki i zamienia pozostałe dwa.
Wymaga to jednego wspólnego koloru dla zamienianej pary,
a kolory dwóch ustalanych wierzchołków można wybrać swobodnie:
każde takie odbicie ustala $2^3=8$ kolorowań.
\item Każde z dwóch odbić względem osi przechodzących przez
środki przeciwległych boków zamienia dwie pary wierzchołków.
Ustala $2^2=4$ kolorowania.
\end{enumerate}
W każdym przypadku liczba kolorowań jest potęgą $2$
o wykładniku równym liczbie cykli odpowiedniej permutacji.
Łączna liczba punktów stałych wynosi
$16+2+2+4+8+8+4+4=48$.
Na mocy lematu Burnside'a liczba różnych kolorowań z dokładnością
do symetrii wynosi $48/8=6$.
Można to sprawdzić, klasyfikując kolorowania według liczby czarnych
wierzchołków: $0$ i $4$ dają po jednej klasie, $1$ i $3$ także
po jednej, a dla dokładnie dwóch czarnych wierzchołków
są dwa rodzaje (sąsiednie lub przeciwległe). Razem sześć.
\end{przyklad}

\begin{uwaga}[Konkretni reprezentanci sześciu orbit]
W obliczonym przykładzie sześć orbit można reprezentować przez
kolorowania z odpowiednio: zerem czarnych wierzchołków,
jednym czarnym, dwoma sąsiednimi czarnymi,
dwoma naprzeciwległymi czarnymi, trzema czarnymi i czterema
czarnymi wierzchołkami. Odbicia i obroty nie zmieniają
liczby czarnych wierzchołków, a w przypadku dokładnie dwóch
czarnych wierzchołków zachowują również ich sąsiedztwo.
Dlatego sześć wymienionych przypadków jest parami nierównoważnych.
Z drugiej strony symetrie kwadratu przeprowadzają dowolne
kolorowanie danego rodzaju na wybranego reprezentanta,
co daje niezależną kontrolę wyniku lematu Burnside'a.
\end{uwaga}

\begin{uwaga}[Kolejny krok: twierdzenia Sylowa]
W całym rozdziale wykorzystaliśmy działania grup na punktach,
warstwach, podgrupach i samych grupach.
W następnym rozdziale zastosujemy je do podgrup o rzędach
będących potęgami liczb pierwszych. W ten sposób wyprowadzimy
klasyczne twierdzenia Cauchy'ego i Sylowa.
\end{uwaga}
